% TOP_PROBLEM: {"id": 3000065, "title": "Partitioning a bipartite graph into proportional factors", "category_id": 2, "category": {"id": 2, "name": "combinatorics", "display_name": "Combinatorics", "description": "Counting problems, graph theory, discrete structures.", "slug": "combinatorics", "order_index": 2, "created_at": "2026-07-31T15:26:25.671Z"}, "status": "open", "problem_number": "AMR-029-0065", "set_id": 13}
% FIRST_POSTED: 2026-10-01
% ENTRY_KIND: statement_only
% UnsolvedMath ID 3000065; source catalogue code AMR-029-0065
% !TeX program = lualatex
% Definition and sources only; this file is not an LLM solution attempt.
\documentclass[11pt,a4paper]{article}
\usepackage[margin=25mm]{geometry}
\usepackage{fontspec}
\setmainfont{lmroman10-regular.otf}[BoldFont=lmroman10-bold.otf,
 ItalicFont=lmroman10-italic.otf,BoldItalicFont=lmroman10-bolditalic.otf]
\usepackage{amsmath,amssymb,mathtools}
\usepackage{unicode-math}
\setmathfont{latinmodern-math.otf}
\AtBeginDocument{\let\setminus\smallsetminus}
\usepackage{xurl,hyperref}
\hypersetup{colorlinks=true,urlcolor=blue}
\setcounter{secnumdepth}{0}
\setlength{\parindent}{0pt}
\setlength{\parskip}{5pt}
\setlength{\emergencystretch}{3em}
\begin{document}
\section*{Partitioning a bipartite graph into proportional factors}\label{3000065}
{\small UnsolvedMath ID 3000065 · AMR-029-0065 · Combinatorics · AMR Open Problem Lists}
\subsection{Definitions and mathematical statement}
Let $G=(U\sqcup W,E)$ be a finite bipartite graph and let
$c_1,\ldots,c_k>0$ be real numbers with $\sum_i c_i=1$.
For $F\subseteq E$, write $d_F(v)$ for the number of edges of $F$
incident with $v$.

Does there always exist a partition $E=E_1\sqcup\cdots\sqcup E_k$
such that
\[
                 \lfloor c_i d_E(v)\rfloor\le d_{E_i}(v)
                       \le\lceil c_i d_E(v)\rceil
                              \quad(v\in U\sqcup W,\ 1\le i\le k)?
\]
Parts are permitted to be empty. The partition specifies spanning
subgraphs on the original vertex set, each approximating a prescribed
fraction of every local degree simultaneously.

Whenever $c_i d_E(v)$ is an integer, exact equality is required;
otherwise the permitted values are its two neighboring integers.
The proportions can differ and do not depend on the vertex.
An edge's part contributes to the degree requirement at both its
endpoints, so independently rounding the desired degrees at each
vertex need not produce a consistent edge partition.
The assertion asks only for this degree control, not for connectivity
or regularity of the parts. Handling equal proportions or just two
parts covers special cases; the general question concerns arbitrary
positive proportions and every finite number of parts.
\subsection{Short English statement}
Can a bipartite graph's edges always be divided into prescribed
proportions while rounding each prescribed vertex degree either down or up to an integer?
\subsection{Sources}
\begin{itemize}
\item[S1] ULAM.AI and the UnsolvedMath Contributors, supplied problems.json, record 3000065 (AMR-029-0065), AMR Open Problem Lists. Catalogue identity and source transcription; CC BY 4.0.
\url{https://huggingface.co/datasets/ulamai/UnsolvedMath}.
\item[S2] Egres Open - List of open problems, Open-problem page: Partitioning a bipartite graph into proportional factors. Original source indexed by AMR-029-0065.
\url{https://oldlemon.cs.elte.hu/egres/open/List_of_open_problems}.
\item[S3] EGRES Open, Partitioning a bipartite graph into proportional factors. Individual source page and explanatory remarks.
\url{https://lemon.cs.elte.hu/egres/open/Partitioning_a_bipartite_graph_into_proportional_factors}.
\end{itemize}
\end{document}
